%%%%%%%%%%%%%%%%%%%%%%%%%%%%%%%%%%%%%%%%%%%%%%%%%%%%%%%%%%%%%%%%%%%
% Header: General study / work / project template
%%%%%%%%%%%%%%%%%%%%%%%%%%%%%%%%%%%%%%%%%%%%%%%%%%%%%%%%%%%%%%%%%%%

\documentclass[11pt,a4paper]{scrartcl}

% Layout
\usepackage[margin=2.5cm]{geometry}
\usepackage[onehalfspacing]{setspace}

% General packages
\usepackage{graphicx}
\usepackage{enumitem}
\usepackage[most]{tcolorbox}
\tcbuselibrary{breakable}
\usepackage[breaklinks=true,colorlinks=true,linkcolor=blue,urlcolor=blue,citecolor=blue]{hyperref}
\usepackage{amsmath,amsthm,amssymb,mathtools}
\usepackage{braket}
\usepackage{icomma}
\usepackage[T1]{fontenc}
\usepackage[utf8]{inputenc}

% Language
\usepackage[ngerman]{babel}

% Units / tables / floats / references
\usepackage[locale = DE,space-before-unit=true,per-mode = symbol]{siunitx}
\usepackage{booktabs,multirow}
\usepackage{placeins}
\usepackage[natbib,abbreviate=true,doi=false,style=numeric-comp,giveninits=true,sorting=none]{biblatex}
\usepackage{csquotes}

% Optional bibliography
% \addbibresource{references.bib}

% Graphics paths
\graphicspath{{Bilder/} {images/} {figures/} {chapters/}}

% Optional math shortcuts
\newcommand{\R}{\mathbb{R}}
\newcommand{\N}{\mathbb{N}}
\newcommand{\Q}{\mathbb{Q}}
\newcommand{\abs}[1]{\left| #1 \right|}
\newcommand{\norm}[1]{\left\lVert #1 \right\rVert}
\newcommand{\ol}[1]{\overline{#1}}
\newcommand{\dd}{\mathrm{d}}
\newcommand{\e}{\mathrm{e}}
\newcommand{\ii}{\mathrm{i}}
\newcommand{\op}[1]{\hat{#1}}
\newcommand{\Ezeroone}{E_1^{(0)}}
\newcommand{\Ezerotwo}{E_2^{(0)}}
\newcommand{\DeltaE}{\Delta}

% Optional units
\DeclareSIUnit{\dBm}{dBm}

% Box styles
\newtcolorbox{calculationbox}[1][]{
    title=Calculation,
    breakable,
    enhanced,
    colback=white,
    colframe=black!60,
    #1
}

\newtcolorbox{solutionbox}[1][]{
    title=Solution,
    breakable,
    enhanced,
    colback=white,
    colframe=black!60,
    #1
}

\newtcolorbox{answerbox}[1][]{
    title=Answer,
    breakable,
    enhanced,
    colback=white,
    colframe=black!60,
    #1
}

\newtcolorbox{notebox}[1][]{
    title=Notes,
    breakable,
    enhanced,
    colback=white,
    colframe=black!60,
    #1
}

\newtcolorbox{sidecalcbox}[1][]{
    title=Side Calculation,
    breakable,
    enhanced,
    colback=white,
    colframe=black!60,
    #1
}

%%%%%%%%%%%%%%%%%%%%%%%%%%%%%%%%%%%%%%%%%%%%%%%%%%%%%%%%%%%%%%%%%%%
\begin{document}

    \title{Theoretische Physik 2 -- Blatt 09}
    \author{Tobias Laser}
    \date{\today}
    \maketitle
    \vfill
    \thispagestyle{empty}
    \newpage

    \pagenumbering{arabic}

    % ================================================================
    \paragraph{21. Zeitunabhängige Störungstheorie und ein 2-Niveau-System:}

    Wir betrachten ein 2-Niveau-System mit Hamiltonoperator
    \begin{align*}
        \op H
        =
        \op H_0 + \lambda \op V
        =
        \begin{pmatrix}
            E_1^{(0)} & 0 \\
            0 & E_2^{(0)}
        \end{pmatrix}
        +
        \lambda
        \begin{pmatrix}
            1 & 1 \\
            1 & 1
        \end{pmatrix}.
    \end{align*}

    \begin{enumerate}[label=(\alph*)]
        \item Finden Sie die Eigenzustände von $\op H_0$ (ohne Rechnung), sowie eine exakte Lösung für die Eigenwerte $E_1,E_2$ und Eigenzustände $\ket{\psi_1},\ket{\psi_2}$ des gesamten Hamiltonoperators $\op H$.

        \begin{calculationbox}
            \begin{solutionbox}
                Zuerst benennen wir die ungestörte Basis. Da $\op H_0$ bereits diagonal gegeben ist, sind die Eigenzustände ohne Rechnung einfach die beiden Standardbasiszustände
                \begin{align*}
                    \ket{1}
                    =
                    \begin{pmatrix}1\\0\end{pmatrix},
                    \qquad
                    \ket{2}
                    =
                    \begin{pmatrix}0\\1\end{pmatrix},
                \end{align*}
                mit
                \begin{align*}
                    \op H_0\ket{1}=E_1^{(0)}\ket{1},
                    \qquad
                    \op H_0\ket{2}=E_2^{(0)}\ket{2}.
                \end{align*}
                Für den gesamten Hamiltonoperator gilt
                \begin{align*}
                    \op H
                    =
                    \begin{pmatrix}
                        E_1^{(0)} + \lambda & \lambda\\
                        \lambda & E_2^{(0)} + \lambda
                    \end{pmatrix}.
                \end{align*}
            \end{solutionbox}

            \begin{calculationbox}[title=Exakte Eigenwerte]
                Die Eigenwerte erhält man aus der charakteristischen Gleichung
                \begin{align*}
                    0
                    &=
                    \det(\op H-E\mathbb{1}) \\
                    &=
                    \det
                    \begin{pmatrix}
                        E_1^{(0)}+\lambda-E & \lambda\\
                        \lambda & E_2^{(0)}+\lambda-E
                    \end{pmatrix} \\
                    &=
                    (E_1^{(0)}+\lambda-E)(E_2^{(0)}+\lambda-E)-\lambda^2.
                \end{align*}
                Eine symmetrische Schreibweise bekommt man mit
                \begin{align*}
                    \bar E^{(0)}
                    =
                    \frac{E_1^{(0)}+E_2^{(0)}}{2},
                    \qquad
                    \Delta
                    =
                    E_1^{(0)}-E_2^{(0)}.
                \end{align*}
                Dann sind die exakten Eigenwerte
                \begin{align*}
                    \boxed{
                    E_{\pm}
                    =
                    \bar E^{(0)}+\lambda
                    \pm
                    \sqrt{\left(\frac{\Delta}{2}\right)^2+\lambda^2}}
                    .
                \end{align*}
                Hier bezeichnet $E_+$ den höheren und $E_-$ den niedrigeren Eigenwert.
            \end{calculationbox}

            \begin{calculationbox}[title=Exakte Eigenzustände]
                Für einen Eigenzustand
                \begin{align*}
                    \ket{\psi_\pm}
                    =
                    a_\pm\ket{1}+b_\pm\ket{2}
                \end{align*}
                gilt
                \begin{align*}
                    (E_1^{(0)}+\lambda-E_\pm)a_\pm
                    +
                    \lambda b_\pm
                    =0.
                \end{align*}
                Daher kann man zum Beispiel
                \begin{align*}
                    b_\pm
                    =
                    \frac{E_\pm-E_1^{(0)}-\lambda}{\lambda}a_\pm
                \end{align*}
                wählen. Mit
                \begin{align*}
                    q_\pm
                    =
                    \frac{E_\pm-E_1^{(0)}-\lambda}{\lambda}
                    =
                    \frac{-\Delta/2 \pm \sqrt{(\Delta/2)^2+\lambda^2}}{\lambda}
                \end{align*}
                ergibt sich die normierte Form
                \begin{align*}
                    \boxed{
                    \ket{\psi_\pm}
                    =
                    \frac{1}{\sqrt{1+q_\pm^2}}
                    \left(\ket{1}+q_\pm\ket{2}\right)}.
                \end{align*}
                Alternativ kann man dieselben Zustände über einen Mischungswinkel $\theta$ schreiben:
                \begin{align*}
                    \tan(2\theta)
                    =
                    \frac{2\lambda}{E_1^{(0)}-E_2^{(0)}}.
                \end{align*}
                Für $E_1^{(0)}>E_2^{(0)}$ kann man dann wählen
                \begin{align*}
                    \ket{\psi_+}=\cos\theta\ket{1}+\sin\theta\ket{2},
                    \qquad
                    \ket{\psi_-}=-\sin\theta\ket{1}+\cos\theta\ket{2}.
                \end{align*}
            \end{calculationbox}
        \end{calculationbox}

        \item Nehmen Sie an, dass $\lambda \ll \abs{E_1^{(0)}-E_2^{(0)}}$, und berechnen Sie mit zeitunabhängiger Störungstheorie die Energieeigenwerte in erster und zweiter Ordnung in $\lambda$, sowie die Energieeigenzustände bis zu erster Ordnung in $\lambda$. Vergleichen Sie dieses Ergebnis mit (a).

        \begin{calculationbox}
            \begin{solutionbox}
                In diesem Teil ist der Abstand der ungestörten Energien groß im Vergleich zur Störung. Daher verwenden wir nicht-entartete Störungstheorie. Wir setzen
                \begin{align*}
                    \Delta = E_1^{(0)}-E_2^{(0)}.
                \end{align*}
                Die Störmatrix hat in der Basis $\{\ket{1},\ket{2}\}$ die Matrixelemente
                \begin{align*}
                    V_{11}=1,
                    \qquad
                    V_{22}=1,
                    \qquad
                    V_{12}=V_{21}=1.
                \end{align*}
            \end{solutionbox}

            \begin{calculationbox}[title=Energie von Zustand $\ket{1}$]
                Für den Zustand $\ket{1}$ ist die Energie in erster Ordnung
                \begin{align*}
                    E_1^{(1)}
                    =
                    \lambda V_{11}
                    =
                    \lambda.
                \end{align*}
                Die Energie in zweiter Ordnung ist
                \begin{align*}
                    E_1^{(2)}
                    =
                    \lambda^2
                    \sum_{m\neq 1}
                    \frac{\abs{V_{m1}}^2}{E_1^{(0)}-E_m^{(0)}}
                    =
                    \lambda^2
                    \frac{\abs{V_{21}}^2}{E_1^{(0)}-E_2^{(0)}}
                    =
                    \frac{\lambda^2}{\Delta}.
                \end{align*}
                Also
                \begin{align*}
                    \boxed{
                    E_1
                    =
                    E_1^{(0)}+\lambda+\frac{\lambda^2}{\Delta}
                    +\mathcal O(\lambda^3)}.
                \end{align*}
            \end{calculationbox}

            \begin{calculationbox}[title=Energie von Zustand $\ket{2}$]
                Analog gilt
                \begin{align*}
                    E_2^{(1)}
                    =
                    \lambda V_{22}
                    =
                    \lambda,
                \end{align*}
                und
                \begin{align*}
                    E_2^{(2)}
                    =
                    \lambda^2
                    \frac{\abs{V_{12}}^2}{E_2^{(0)}-E_1^{(0)}}
                    =
                    -\frac{\lambda^2}{\Delta}.
                \end{align*}
                Also
                \begin{align*}
                    \boxed{
                    E_2
                    =
                    E_2^{(0)}+\lambda-\frac{\lambda^2}{\Delta}
                    +\mathcal O(\lambda^3)}.
                \end{align*}
            \end{calculationbox}

            \begin{calculationbox}[title=Eigenzustände bis erster Ordnung]
                Für die Korrektur der Eigenzustände gilt in erster Ordnung
                \begin{align*}
                    \ket{n^{(1)}}
                    =
                    \lambda
                    \sum_{m\neq n}
                    \frac{V_{mn}}{E_n^{(0)}-E_m^{(0)}}\ket{m}.
                \end{align*}
                Damit folgt
                \begin{align*}
                    \boxed{
                    \ket{\psi_1}
                    =
                    \ket{1}
                    +
                    \frac{\lambda}{\Delta}\ket{2}
                    +\mathcal O(\lambda^2)}
                \end{align*}
                und
                \begin{align*}
                    \boxed{
                    \ket{\psi_2}
                    =
                    \ket{2}
                    -
                    \frac{\lambda}{\Delta}\ket{1}
                    +\mathcal O(\lambda^2)}.
                \end{align*}
                Die Normierung unterscheidet sich erst in Ordnung $\lambda^2$, daher ist diese Schreibweise bis erster Ordnung ausreichend.
            \end{calculationbox}

            \begin{solutionbox}
                Der Vergleich mit der exakten Lösung aus (a) funktioniert durch Entwicklung der Wurzel für $\abs{\lambda}\ll\abs{\Delta}$. Für $\Delta>0$ gilt zum Beispiel
                \begin{align*}
                    \sqrt{\left(\frac{\Delta}{2}\right)^2+\lambda^2}
                    =
                    \frac{\Delta}{2}
                    \sqrt{1+\frac{4\lambda^2}{\Delta^2}}
                    =
                    \frac{\Delta}{2}
                    +
                    \frac{\lambda^2}{\Delta}
                    +
                    \mathcal O(\lambda^4).
                \end{align*}
                Einsetzen liefert genau die Energien aus der Störungstheorie. Auch die Eigenzustände stimmen überein, wenn man die exakten Koeffizienten $q_\pm$ bis zur ersten Ordnung in $\lambda$ entwickelt.
            \end{solutionbox}
        \end{calculationbox}

        \item Nehmen Sie an, dass $E_1^{(0)}$ und $E_2^{(0)}$ fast entartet sind, d.h. $\abs{E_1^{(0)}-E_2^{(0)}}\ll \lambda$, und berechnen Sie die Energieeigenwerte bis zu erster Ordnung in $\lambda$, sowie die Energieeigenzustände nullter Ordnung mit Störungstheorie im Fall $E_1^{(0)}=E_2^{(0)}$. Vergleichen Sie das Ergebnis mit der exakten Lösung von (a).

        \begin{calculationbox}
            \begin{solutionbox}
                Wenn die ungestörten Energien entartet oder fast entartet sind, darf man nicht-entartete Störungstheorie nicht direkt verwenden, weil die Nenner $E_n^{(0)}-E_m^{(0)}$ klein werden. Deshalb betrachten wir zuerst den exakt entarteten Fall
                \begin{align*}
                    E_1^{(0)}=E_2^{(0)}=:E^{(0)}.
                \end{align*}
                Dann ist die Störmatrix im entarteten Unterraum
                \begin{align*}
                    \lambda V
                    =
                    \lambda
                    \begin{pmatrix}
                        1&1\\
                        1&1
                    \end{pmatrix}.
                \end{align*}
            \end{solutionbox}

            \begin{calculationbox}[title=Diagonalisierung der Störung im entarteten Unterraum]
                Wir diagonalisieren die Matrix
                \begin{align*}
                    V
                    =
                    \begin{pmatrix}
                        1&1\\
                        1&1
                    \end{pmatrix}.
                \end{align*}
                Für den symmetrischen Zustand gilt
                \begin{align*}
                    V\frac{1}{\sqrt2}
                    \begin{pmatrix}1\\1\end{pmatrix}
                    =
                    \frac{1}{\sqrt2}
                    \begin{pmatrix}2\\2\end{pmatrix}
                    =
                    2\frac{1}{\sqrt2}
                    \begin{pmatrix}1\\1\end{pmatrix}.
                \end{align*}
                Für den antisymmetrischen Zustand gilt
                \begin{align*}
                    V\frac{1}{\sqrt2}
                    \begin{pmatrix}1\\-1\end{pmatrix}
                    =
                    \frac{1}{\sqrt2}
                    \begin{pmatrix}0\\0\end{pmatrix}.
                \end{align*}
                Damit sind die richtigen Zustände nullter Ordnung
                \begin{align*}
                    \boxed{
                    \ket{\psi_+^{(0)}}
                    =
                    \frac{1}{\sqrt2}(\ket{1}+\ket{2})}
                \end{align*}
                und
                \begin{align*}
                    \boxed{
                    \ket{\psi_-^{(0)}}
                    =
                    \frac{1}{\sqrt2}(\ket{1}-\ket{2})}.
                \end{align*}
            \end{calculationbox}

            \begin{calculationbox}[title=Energien bis erster Ordnung]
                Die Energieverschiebungen erster Ordnung sind die Eigenwerte der Störmatrix $\lambda V$. Daher gilt
                \begin{align*}
                    \boxed{
                    E_+
                    =
                    E^{(0)}+2\lambda}
                \end{align*}
                und
                \begin{align*}
                    \boxed{
                    E_-
                    =
                    E^{(0)}}.
                \end{align*}
            \end{calculationbox}

            \begin{solutionbox}
                Vergleicht man mit der exakten Lösung aus (a), setzt man dort $\Delta=0$. Dann folgt
                \begin{align*}
                    E_\pm
                    =
                    E^{(0)}+\lambda\pm\abs{\lambda}.
                \end{align*}
                Für $\lambda>0$ ist also
                \begin{align*}
                    E_+=E^{(0)}+2\lambda,
                    \qquad
                    E_-=E^{(0)}.
                \end{align*}
                Genau das ist das Ergebnis der entarteten Störungstheorie. Der wichtige Punkt ist: Bei Entartung muss man zuerst die Störung im entarteten Unterraum diagonalisieren. Die ungestörten Zustände $\ket{1}$ und $\ket{2}$ sind dann nicht mehr die geeignete Basis.
            \end{solutionbox}
        \end{calculationbox}
    \end{enumerate}

    % ================================================================
    \paragraph{22. Magnetische Dipolwechselwirkung:}

    Wir betrachten zwei Spin-$\frac12$-Teilchen mit magnetischem Moment
    \begin{align*}
        \vec M^{(i)}=\gamma \vec S^{(i)},
        \qquad
        i=1,2,
    \end{align*}
    in einem externen Magnetfeld
    \begin{align*}
        \vec B=B_0\vec e_z.
    \end{align*}
    Der Hamiltonoperator des Systems ist damit
    \begin{align*}
        \op H_0
        =
        -\vec B\cdot(\vec M^{(1)}+\vec M^{(2)}).
    \end{align*}

    \begin{enumerate}[label=(\alph*)]
        \item Zeigen Sie, dass der Hamiltonoperator $\op H_0$ diagonal in der Produktbasis
        \begin{align*}
            \ket{++},\quad \ket{+-},\quad \ket{-+},\quad \ket{--}
        \end{align*}
        ist und geben Sie die zugehörigen Energieeigenwerte an. Gibt es entartete Energieniveaus? Es ist nützlich, die Frequenz $\omega_0=\gamma B_0$ zu definieren.

        \begin{calculationbox}
            \begin{solutionbox}
                Da das Magnetfeld in $z$-Richtung zeigt, kommt nur die $z$-Komponente der Spins vor. Mit $\omega_0=\gamma B_0$ ist
                \begin{align*}
                    \op H_0
                    =
                    -B_0\gamma\left(\op S_z^{(1)}+\op S_z^{(2)}\right)
                    =
                    -\omega_0\left(\op S_z^{(1)}+\op S_z^{(2)}\right).
                \end{align*}
                Die Produktzustände sind Eigenzustände von $\op S_z^{(1)}$ und $\op S_z^{(2)}$. Deshalb ist $\op H_0$ in dieser Basis diagonal.
            \end{solutionbox}

            \begin{calculationbox}[title=Energieeigenwerte von $\op H_0$]
                Für Spin-$\frac12$ gilt
                \begin{align*}
                    \op S_z\ket{+}=\frac{\hbar}{2}\ket{+},
                    \qquad
                    \op S_z\ket{-}=-\frac{\hbar}{2}\ket{-}.
                \end{align*}
                Daher erhält man
                \begin{align*}
                    \op H_0\ket{++}
                    &=
                    -\omega_0\left(\frac{\hbar}{2}+\frac{\hbar}{2}\right)\ket{++}
                    =
                    -\hbar\omega_0\ket{++},\\
                    \op H_0\ket{+-}
                    &=
                    -\omega_0\left(\frac{\hbar}{2}-\frac{\hbar}{2}\right)\ket{+-}
                    =
                    0,\\
                    \op H_0\ket{-+}
                    &=
                    -\omega_0\left(-\frac{\hbar}{2}+\frac{\hbar}{2}\right)\ket{-+}
                    =
                    0,\\
                    \op H_0\ket{--}
                    &=
                    -\omega_0\left(-\frac{\hbar}{2}-\frac{\hbar}{2}\right)\ket{--}
                    =
                    +\hbar\omega_0\ket{--}.
                \end{align*}
                Also
                \begin{align*}
                    \boxed{E_{++}^{(0)}=-\hbar\omega_0,}
                    \qquad
                    \boxed{E_{+-}^{(0)}=0,}
                    \qquad
                    \boxed{E_{-+}^{(0)}=0,}
                    \qquad
                    \boxed{E_{--}^{(0)}=+\hbar\omega_0.}
                \end{align*}
            \end{calculationbox}

            \begin{solutionbox}
                Die Zustände $\ket{+-}$ und $\ket{-+}$ besitzen dieselbe Energie $0$. Dieses Energieniveau ist also zweifach entartet. Die Zustände $\ket{++}$ und $\ket{--}$ sind im Allgemeinen nicht entartet, solange $\omega_0\neq0$ gilt.
            \end{solutionbox}
        \end{calculationbox}

        \item Die Teilchen sind sich so nahe, dass die Dipol-Dipol-Wechselwirkung der beiden magnetischen Momente nicht vernachlässigt werden kann. Diese wird durch einen Hamiltonoperator
        \begin{align*}
            \op H_1
            =
            \frac{\mu_0}{4\pi}\frac{1}{r^3}
            \left[
                \vec M^{(1)}\cdot \vec M^{(2)}
                -3(\vec M^{(1)}\cdot\vec n)(\vec M^{(2)}\cdot\vec n)
            \right]
        \end{align*}
        beschrieben, wobei $r$ den Abstand zwischen den Teilchen und
        \begin{align*}
            \vec n
            =
            (\sin\theta\cos\varphi,\sin\theta\sin\varphi,\cos\theta)
        \end{align*}
        den Einheitsvektor in diese Richtung bezeichnet. Der Abstand $r$ sei groß genug, sodass $\op H_1$ als Störungsterm behandelt werden kann. Berechnen Sie zuerst die Verschiebung der Energieniveaus $\ket{++}$, $\ket{--}$ in erster Ordnung von $\op H_1$. Zur Berechnung der relevanten Matrixelemente $\bra{\pm\pm}\op H_1\ket{\pm\pm}$ müssen nur die Terme in $\op H_1$ betrachtet werden, die proportional zu $S_z^{(1)}S_z^{(2)}$ sind. Warum? Es ist nützlich, die Frequenz
        \begin{align*}
            \Omega
            =
            \frac{\mu_0\gamma^2\hbar}{4\pi r^3}
        \end{align*}
        zu definieren.

        \begin{calculationbox}
            \begin{solutionbox}
                Wir schreiben zuerst $\op H_1$ mit $\vec M^{(i)}=\gamma\vec S^{(i)}$. Dann ist
                \begin{align*}
                    \op H_1
                    =
                    \frac{\mu_0\gamma^2}{4\pi r^3}
                    \left[
                        \vec S^{(1)}\cdot \vec S^{(2)}
                        -3(\vec S^{(1)}\cdot\vec n)(\vec S^{(2)}\cdot\vec n)
                    \right].
                \end{align*}
                Wegen
                \begin{align*}
                    \Omega
                    =
                    \frac{\mu_0\gamma^2\hbar}{4\pi r^3}
                \end{align*}
                gilt
                \begin{align*}
                    \frac{\mu_0\gamma^2}{4\pi r^3}
                    =
                    \frac{\Omega}{\hbar}.
                \end{align*}
            \end{solutionbox}

            \begin{calculationbox}[title=Warum nur $S_z^{(1)}S_z^{(2)}$?]
                Die Zustände $\ket{++}$ und $\ket{--}$ sind Eigenzustände von $S_z^{(1)}$ und $S_z^{(2)}$. Für ein diagonales Matrixelement wie $\bra{++}\op H_1\ket{++}$ müssen Anfangs- und Endzustand gleich sein. Operatoren $S_x$ und $S_y$ enthalten aber Leiteroperatoren und ändern den Spin:
                \begin{align*}
                    S_x=\frac12(S_++S_-),
                    \qquad
                    S_y=\frac{1}{2\ii}(S_+-S_-).
                \end{align*}
                Dadurch entstehen Zustände, die orthogonal zu $\ket{++}$ beziehungsweise $\ket{--}$ sind. Deshalb tragen diese Terme zu den diagonalen Erwartungswerten nicht bei. Nur der Anteil proportional zu
                \begin{align*}
                    S_z^{(1)}S_z^{(2)}
                \end{align*}
                bleibt übrig.
            \end{calculationbox}

            \begin{calculationbox}[title=Koeffizient vor $S_z^{(1)}S_z^{(2)}$]
                Aus
                \begin{align*}
                    \vec S^{(1)}\cdot\vec S^{(2)}
                    &=
                    S_x^{(1)}S_x^{(2)}+S_y^{(1)}S_y^{(2)}+S_z^{(1)}S_z^{(2)},\\
                    (\vec S^{(1)}\cdot\vec n)(\vec S^{(2)}\cdot\vec n)
                    &\supset
                    \cos^2\theta\,S_z^{(1)}S_z^{(2)}
                \end{align*}
                folgt für den $S_z^{(1)}S_z^{(2)}$-Anteil
                \begin{align*}
                    S_z^{(1)}S_z^{(2)}
                    -3\cos^2\theta\,S_z^{(1)}S_z^{(2)}
                    =
                    (1-3\cos^2\theta)S_z^{(1)}S_z^{(2)}.
                \end{align*}
            \end{calculationbox}

            \begin{calculationbox}[title=Energieverschiebung für $\ket{++}$ und $\ket{--}$]
                Für beide Zustände gilt
                \begin{align*}
                    \bra{++}S_z^{(1)}S_z^{(2)}\ket{++}
                    &=
                    \frac{\hbar}{2}\frac{\hbar}{2}
                    =
                    \frac{\hbar^2}{4},\\
                    \bra{--}S_z^{(1)}S_z^{(2)}\ket{--}
                    &=
                    \left(-\frac{\hbar}{2}\right)
                    \left(-\frac{\hbar}{2}\right)
                    =
                    \frac{\hbar^2}{4}.
                \end{align*}
                Daher ist die Verschiebung erster Ordnung
                \begin{align*}
                    \Delta E_{++}^{(1)}
                    =
                    \Delta E_{--}^{(1)}
                    &=
                    \frac{\Omega}{\hbar}
                    (1-3\cos^2\theta)
                    \frac{\hbar^2}{4}\\
                    &=
                    \boxed{
                    \frac{\hbar\Omega}{4}(1-3\cos^2\theta)}.
                \end{align*}
            \end{calculationbox}
        \end{calculationbox}

        \item Berechnen Sie nun die Verschiebung der Energieniveaus $\ket{+-}$, $\ket{-+}$ in erster Ordnung von $\op H_1$ und bestimmen Sie die zugehörigen Spinwellenfunktionen. Zur Berechnung der relevanten Matrixelemente $\bra{\pm\mp}\op H_1\ket{\pm\mp}$ bzw. $\bra{\pm\mp}\op H_1\ket{\mp\pm}$ müssen nur die Terme in $\op H_1$ betrachtet werden, die proportional zu $S_z^{(1)}S_z^{(2)}$ bzw. $S_x^{(1)}S_x^{(2)}$, $S_y^{(1)}S_y^{(2)}$ sind. Warum?

        \begin{calculationbox}
            \begin{solutionbox}
                Die Zustände $\ket{+-}$ und $\ket{-+}$ sind bezüglich $\op H_0$ entartet. Deshalb müssen wir die Störung $\op H_1$ im zweidimensionalen Unterraum
                \begin{align*}
                    \mathcal H_{0}
                    =
                    \operatorname{span}\{\ket{+-},\ket{-+}\}
                \end{align*}
                diagonalisieren. Das ist entartete Störungstheorie.
            \end{solutionbox}

            \begin{calculationbox}[title=Warum diese Terme relevant sind]
                Für die diagonalen Matrixelemente
                \begin{align*}
                    \bra{+-}\op H_1\ket{+-},
                    \qquad
                    \bra{-+}\op H_1\ket{-+}
                \end{align*}
                muss der Zustand unverändert bleiben. Dafür trägt wieder nur $S_z^{(1)}S_z^{(2)}$ bei.

                Für die außerdiagonalen Matrixelemente
                \begin{align*}
                    \bra{+-}\op H_1\ket{-+},
                    \qquad
                    \bra{-+}\op H_1\ket{+-}
                \end{align*}
                müssen beide Spins gleichzeitig umgedreht werden. Genau das können die Terme $S_x^{(1)}S_x^{(2)}$ und $S_y^{(1)}S_y^{(2)}$ leisten. Terme mit nur einem $S_z$ und einem $S_x$ oder $S_y$ ändern nur einen Spin oder liefern wegen Orthogonalität keinen Beitrag. Die gemischten $xy$- und $yx$-Terme heben sich in der Summe der beiden Matrixelemente auf, weil ihre Beiträge rein imaginär mit entgegengesetztem Vorzeichen sind.
            \end{calculationbox}

            \begin{calculationbox}[title=Diagonale Matrixelemente]
                Der diagonale Anteil ist wie zuvor
                \begin{align*}
                    \frac{\Omega}{\hbar}
                    (1-3\cos^2\theta)S_z^{(1)}S_z^{(2)}.
                \end{align*}
                Für $\ket{+-}$ gilt
                \begin{align*}
                    \bra{+-}S_z^{(1)}S_z^{(2)}\ket{+-}
                    =
                    \frac{\hbar}{2}\left(-\frac{\hbar}{2}\right)
                    =
                    -\frac{\hbar^2}{4}.
                \end{align*}
                Dasselbe gilt für $\ket{-+}$. Also
                \begin{align*}
                    \bra{+-}\op H_1\ket{+-}
                    =
                    \bra{-+}\op H_1\ket{-+}
                    =
                    -\frac{\hbar\Omega}{4}(1-3\cos^2\theta).
                \end{align*}
                Wir schreiben dafür
                \begin{align*}
                    B
                    :=
                    \frac{\hbar\Omega}{4}(3\cos^2\theta-1)
                    =
                    -\frac{\hbar\Omega}{4}(1-3\cos^2\theta).
                \end{align*}
                Dann sind die diagonalen Matrixelemente gleich $B$.
            \end{calculationbox}

            \begin{calculationbox}[title=Außerdiagonale Matrixelemente]
                Für die außerdiagonalen Matrixelemente brauchen wir die Beiträge von $S_x^{(1)}S_x^{(2)}$ und $S_y^{(1)}S_y^{(2)}$. Die zugehörigen Koeffizienten in $\op H_1$ sind
                \begin{align*}
                    S_x^{(1)}S_x^{(2)} &: \quad 1-3\sin^2\theta\cos^2\varphi,\\
                    S_y^{(1)}S_y^{(2)} &: \quad 1-3\sin^2\theta\sin^2\varphi.
                \end{align*}
                Außerdem gilt
                \begin{align*}
                    \bra{+-}S_x^{(1)}S_x^{(2)}\ket{-+}
                    &=
                    \frac{\hbar^2}{4},\\
                    \bra{+-}S_y^{(1)}S_y^{(2)}\ket{-+}
                    &=
                    \frac{\hbar^2}{4}.
                \end{align*}
                Damit folgt
                \begin{align*}
                    \bra{+-}\op H_1\ket{-+}
                    &=
                    \frac{\Omega}{\hbar}
                    \frac{\hbar^2}{4}
                    \left[
                        1-3\sin^2\theta\cos^2\varphi
                        +1-3\sin^2\theta\sin^2\varphi
                    \right]\\
                    &=
                    \frac{\hbar\Omega}{4}
                    \left[2-3\sin^2\theta\right]\\
                    &=
                    \frac{\hbar\Omega}{4}
                    \left[3\cos^2\theta-1\right]\\
                    &=
                    B.
                \end{align*}
                Also ist auch
                \begin{align*}
                    \bra{-+}\op H_1\ket{+-}=B.
                \end{align*}
            \end{calculationbox}

            \begin{calculationbox}[title=Störmatrix und Diagonalisierung]
                Im entarteten Unterraum $\{\ket{+-},\ket{-+}\}$ lautet die Störmatrix erster Ordnung
                \begin{align*}
                    H_1^{\mathrm{eff}}
                    =
                    \begin{pmatrix}
                        B&B\\
                        B&B
                    \end{pmatrix},
                    \qquad
                    B
                    =
                    \frac{\hbar\Omega}{4}(3\cos^2\theta-1).
                \end{align*}
                Diese Matrix hat die Eigenvektoren
                \begin{align*}
                    \frac{1}{\sqrt2}
                    \begin{pmatrix}1\\1\end{pmatrix},
                    \qquad
                    \frac{1}{\sqrt2}
                    \begin{pmatrix}1\\-1\end{pmatrix},
                \end{align*}
                mit Eigenwerten
                \begin{align*}
                    2B,
                    \qquad
                    0.
                \end{align*}
                Daher sind die richtigen Spinwellenfunktionen
                \begin{align*}
                    \boxed{
                    \ket{\psi_s}
                    =
                    \frac{1}{\sqrt2}(\ket{+-}+\ket{-+})}
                \end{align*}
                mit Energieverschiebung
                \begin{align*}
                    \boxed{
                    \Delta E_s^{(1)}
                    =
                    2B
                    =
                    \frac{\hbar\Omega}{2}(3\cos^2\theta-1)}
                \end{align*}
                und
                \begin{align*}
                    \boxed{
                    \ket{\psi_a}
                    =
                    \frac{1}{\sqrt2}(\ket{+-}-\ket{-+})}
                \end{align*}
                mit Energieverschiebung
                \begin{align*}
                    \boxed{
                    \Delta E_a^{(1)}=0}.
                \end{align*}
            \end{calculationbox}
        \end{calculationbox}

        \item Skizzieren Sie, wie sich das Energiespektrum durch die Wechselwirkung verschiebt.

        \begin{calculationbox}
            \begin{solutionbox}
                Ohne Dipol-Dipol-Wechselwirkung hat das System drei Energieniveaus:
                \begin{align*}
                    E_{++}^{(0)}=-\hbar\omega_0,
                    \qquad
                    E_{+-}^{(0)}=E_{-+}^{(0)}=0,
                    \qquad
                    E_{--}^{(0)}=+\hbar\omega_0.
                \end{align*}
                Die Wechselwirkung verschiebt die nicht-entarteten Zustände $\ket{++}$ und $\ket{--}$ beide um denselben Betrag
                \begin{align*}
                    A
                    :=
                    \frac{\hbar\Omega}{4}(1-3\cos^2\theta).
                \end{align*}
                Das entartete Niveau bei $E=0$ spaltet auf: Der symmetrische Zustand wird um
                \begin{align*}
                    -2A
                    =
                    \frac{\hbar\Omega}{2}(3\cos^2\theta-1)
                \end{align*}
                verschoben, der antisymmetrische Zustand bleibt in erster Ordnung unverändert.
            \end{solutionbox}

            \begin{calculationbox}[title=Zusammenfassung des gestörten Spektrums]
                Bis zur ersten Ordnung in der Dipol-Dipol-Wechselwirkung erhält man
                \begin{align*}
                    \boxed{
                    E_{++}
                    =
                    -\hbar\omega_0
                    +
                    \frac{\hbar\Omega}{4}(1-3\cos^2\theta)}
                \end{align*}
                zum Zustand $\ket{++}$,
                \begin{align*}
                    \boxed{
                    E_{--}
                    =
                    +\hbar\omega_0
                    +
                    \frac{\hbar\Omega}{4}(1-3\cos^2\theta)}
                \end{align*}
                zum Zustand $\ket{--}$,
                \begin{align*}
                    \boxed{
                    E_s
                    =
                    \frac{\hbar\Omega}{2}(3\cos^2\theta-1)}
                \end{align*}
                zum symmetrischen Zustand $\ket{\psi_s}$, und
                \begin{align*}
                    \boxed{
                    E_a
                    =
                    0}
                \end{align*}
                zum antisymmetrischen Zustand $\ket{\psi_a}$.
            \end{calculationbox}

            \begin{solutionbox}
                Für die Skizze zeichnet man zuerst die drei ungestörten Niveaus $-\hbar\omega_0$, $0$ und $+\hbar\omega_0$. Das mittlere Niveau ist doppelt entartet. Danach wird die Dipol-Dipol-Wechselwirkung eingeschaltet: Die äußeren Niveaus verschieben sich beide um $A$, während das mittlere Niveau in zwei Linien aufspaltet, nämlich in $E_s=-2A$ und $E_a=0$. Das Vorzeichen der Verschiebung hängt vom Winkel $\theta$ ab. Beim sogenannten magischen Winkel
                \begin{align*}
                    3\cos^2\theta-1=0
                    \qquad\Longleftrightarrow\qquad
                    \theta=\arccos\left(\frac{1}{\sqrt3}\right)
                \end{align*}
                verschwinden diese ersten Ordnungsverschiebungen.
            \end{solutionbox}
        \end{calculationbox}
    \end{enumerate}

\end{document}
